% ===========================================================
% 《线性代数9讲》第 5 讲  线性方程组 —— 片段 A（本讲开头）
% 源：线代9讲强化.pdf  PDF p60–p64（印刷页 49–53），共 5 页
%
% 处理说明：
%   1) 页首「三向解题法」思维导图（PDF p60 = 印刷页 49，占整页）按 AGENTS.md §7.1
%      决定 2 **不转写、不重绘**，只留 `% FIGURE-OPD:` 一行注释。导图中的数学内容
%      「齐次线性方程组 $Ax=0$」「非齐次线性方程组 $Ax=b$」在 PDF p62（印刷页 51）
%      正文中重新出现，已按正文完整转写；
%   2) OPD 方法论标记剔除（依 AGENTS.md §7.1 决定 2）：
%      · \fsec 标题里的 OPD 括号删除，只保留标题文字与路线/方法名：
%        「一、线性方程组理论总结」（$D_1$（常规操作））
%        「1. 齐次线性方程组 $Ax=0$」（$D_1$（常规操作））
%        「2. 非齐次线性方程组 $Ax=b$」（$D_1$（常规操作））
%        「二、线性方程组问题」（$D_1$（常规操作）+$D_{22}$（转换等价表述）+$P_2$（反证思路））
%        「1. 一般求解问题（含参常考）」（$O_1$（盯住目标1）+$D_1$（常规操作）+$P_2$（反证思路））
%        「3. 线性方程组的几何意义（仅数学一）」
%        （$O_4$（盯住目标4）+$D_1$（常规操作）+$D_{22}$（转换等价表述）+$D_{43}$（数形结合））
%      · 正文中指向公式/结论的蓝色纯 OPD 代码与「见到……想到……」方法论语句一律
%        以 `% OPD 蓝色标注（原稿）：……` 注释留痕，不排入正文（同 ch04_b 的处理）：
%        →$D_1$（常规操作），牢记（p62「基础解系」处）
%        →$D_{22}$（转换等价表述），且此条件较强，要留意（p62「解与系数的关系」处）
%        →$D_1$（常规操作），牢记（p62「非齐次线性方程组解的情形」处）
%        →$D_1$（常规操作），牢记（p63【注】(1) 处）
%        $D_{22}$（转换等价表述）；见到 $AB=O$，想到 $AX=O$；见到 $|A|=0$，想到
%          $AA^{*}=O$，$A^{*}A=O$，想到 $AX=O$，$A^{*}X=O$（p63 例5.1 处）
%        →$D_{21}$（观察研究对象），第3讲，矩阵乘法的形状大观（p64 例5.2 处）
%      · 数学符号一律照抄（如 $A$，$\bar A$，$\alpha_i$，$\beta_i$，$D_i$ 作为行列式记号）；
%   3) 蓝色**数学**旁注照原文内联排入（\quad(\text{...}) 或独立成行）：
%      「表达平面 $\Pi_i$ 的方向」「表达确定方向后 $\Pi_i$ 的位置」（p61）；
%      「任何两个面都不重合／存在两个面重合／任何两个面都相交／存在两个面平行但不重合」
%      （p61 两表末列）；「①是 $Ax=0$ 的解；②$\xi_1,\cdots,\xi_{n-r}$ 线性无关；③$s=n-r(A)$」
%      「$Ax=b$ 的一个特解」（p63 通解式旁注）；
%      「$Ax=b$ 的"广义基础解系"」「解中有 $n-r(A)+1$ 个线性无关的向量」（p63【注】(3) 处）；
%   4) 印刷页 50（PDF p61）两张「方程组有解/无解的情形」表：原表三列
%      「图形｜几何意义｜代数表达」。按规范，**图形列不重绘**，留 `% FIGURE:` 占位注释；
%      「几何意义」列作 `\fcard` 标题，「代数表达」列作卡片正文；两表分别用
%      `\fgroup` 标出表头「方程组有解的情形」「方程组无解的情形」；行序照原表自上而下；
%   5) 本片无「习题」栏（已逐页核实 PDF p60–p64）；第 5 讲自印刷页 49 开始；
%   6) 印刷页 51（PDF p62）「1.」「2.」为原稿蓝色粗体子标题（层级在「一、」「二、」之下，
%      在「（1）（2）（3）」之上），本片段用 `\fsec` 承接（可用宏只有两个层级）。
%
% 接缝说明：
%   · 首行承第 4 讲末尾（PDF p59 = 印刷页 48 末）——该讲已结束，本片自思维导图页起；
%   · 例5.1 的解答跨 PDF p63（印刷页 52）末行「……且」与 PDF p64（印刷页 53）首行
%     「$r(B^{\mathrm T})=r[(B^{\mathrm T},b)]=3$……」，本片按一个 `fcard` 排完；
%   · 例5.3 题干在 PDF p64（印刷页 53）末行只到「（1）……；」，其「（2）」小问与两问的
%     【解】在 PDF p65（印刷页 54，属下一分片 ch05_b）；本片在该处结束 fcard。
% ===========================================================

\fchap{第5章\quad 线性方程组}

% -----------------------------------------------------------
% PDF p60（印刷页 49）——整页为「三向解题法」思维导图
% -----------------------------------------------------------
% FIGURE-OPD: 49 三向解题法思维导图（按规则不保留）

% -----------------------------------------------------------
% PDF p61（印刷页 50）
% -----------------------------------------------------------
\fsec{一、线性方程组理论总结}

\fsec{1. 齐次线性方程组 $Ax=0$}

（1）解的性质.

对于齐次线性方程组
\fml{A_{m\times n}x=0,}

若 $\xi_1$，$\xi_2$ 是方程组 $Ax=0$ 的解，则 $x=k_1\xi_1+k_2\xi_2$（$k_1$，$k_2\in\mathbb{R}$）也是它的解.

（2）基础解系与通解.

①设 $r(A)<n$，$\xi_1$，$\xi_2$，$\cdots$，$\xi_s$ 是方程组 $Ax=0$ 的一组解向量. 如果

a. $\xi_1$，$\xi_2$，$\cdots$，$\xi_s$ 线性无关；

b. 方程组 $Ax=0$ 的任一解向量均可由 $\xi_1$，$\xi_2$，$\cdots$，$\xi_s$ 线性表示，即 $s=n-r(A)$，则称 $\xi_1$，$\xi_2$，$\cdots$，$\xi_s$ 是方程组 $Ax=0$ 的一个基础解系.

% OPD 蓝色标注（原稿）：→$D_1$（常规操作），牢记

\begin{fnote}
【注】基础解系应满足三个条件：①是解；②线性无关；③ $s=n-r(A)$.
\end{fnote}

②齐次线性方程组 $Ax=0$ 有非零解的充要条件是 $r(A)<n$，此时它的通解为
\fml{x=k_1\xi_1+k_2\xi_2+\cdots+k_{n-r}\xi_{n-r},}
其中 $r=r(A)$，$k_1$，$k_2$，$\cdots$，$k_{n-r}$ 为任意常数.

（3）解与系数的关系.

若齐次线性方程组
\fml{\begin{cases}a_{11}x_1+a_{12}x_2+\cdots+a_{1n}x_n=0,\\a_{21}x_1+a_{22}x_2+\cdots+a_{2n}x_n=0,\\\cdots\cdots\\a_{m1}x_1+a_{m2}x_2+\cdots+a_{mn}x_n=0\end{cases}}
有解 $\beta=[b_1,b_2,\cdots,b_n]^{\mathrm T}$，即
\fml{a_{i1}b_1+a_{i2}b_2+\cdots+a_{in}b_n=0\quad(i=1,2,\cdots,m).}

% OPD 蓝色标注（原稿）：→$D_{22}$（转换等价表述），且此条件较强，要留意

记系数矩阵 $A$ 的行向量为 $\alpha_i=[a_{i1},a_{i2},\cdots,a_{in}]$（$i=1,2,\cdots,m$），则上式即为
\fml{\alpha_i\beta=0\quad(i=1,2,\cdots,m),}
故系数矩阵 $A$ 的行向量与 $Ax=0$ 的解向量正交.

\fsec{2. 非齐次线性方程组 $Ax=b$}

（1）解的性质.

对于非齐次线性方程组
\fml{A_{m\times n}x=b,}
若 $\eta_1$，$\eta_2$ 都是方程组 $Ax=b$ 的解，则 $\eta_1-\eta_2$ 是相应的齐次线性方程组 $Ax=0$ 的解.

（2）非齐次线性方程组解的情形.

% OPD 蓝色标注（原稿）：→$D_1$（常规操作），牢记

①当 $r(A)=r([A,b])=n$ 时，方程组有唯一解；

% -----------------------------------------------------------
% PDF p63（印刷页 52）
% -----------------------------------------------------------
②当 $r(A)=r([A,b])<n$ 时，方程组有无穷多解；

③当 $r(A)\ne r([A,b])$ 时，方程组无解.

（3）非齐次线性方程组的通解.

设 $r(A)=r<n$，若 $\xi_1$，$\xi_2$，$\cdots$，$\xi_{n-r}$ 为非齐次线性方程组 $Ax=b$ 相应的齐次线性方程组 $Ax=0$ 的基础解系，$\eta^*$ 为非齐次线性方程组 $Ax=b$ 的一个特解，则非齐次线性方程组 $Ax=b$ 的通解为
\fml{x=\eta^*+k_1\xi_1+k_2\xi_2+\cdots+k_{n-r}\xi_{n-r},}

（$\eta^*$ 为 $Ax=b$ 的一个特解；①是 $Ax=0$ 的解；②$\xi_1$，$\cdots$，$\xi_{n-r}$ 线性无关；③ $s=n-r(A)$）

其中 $k_1$，$k_2$，$\cdots$，$k_{n-r}$ 为任意常数.

\begin{fnote}
【注】（1）与非齐次线性方程组 $Ax=b$ 对应的齐次线性方程组 $Ax=0$ 称为非齐次线性方程组 $Ax=b$ 的导出组，故上述结论可简述为非齐次线性方程组的通解等于它的一个特解与其导出组的通解之和.

% OPD 蓝色标注（原稿）：→$D_1$（常规操作），牢记

（2）当未知数个数等于方程个数，且 $|A|\ne0$ 时，可用克拉默法则求解，即 $x_i=\dfrac{D_i}{|A|}$，$i=1,2,\cdots,n$.

（3）设 $\eta^*$ 是非齐次线性方程组 $A_{m\times n}x=b$（$b\ne0$）的一个解，$\xi_1$，$\cdots$，$\xi_{n-r}$ 是其导出组 $Ax=0$ 的一个基础解系，且 $r(A)=r$，则

① $\eta^*$，$\eta^*+\xi_1$，$\cdots$，$\eta^*+\xi_{n-r}$ 线性无关；\quad($Ax=b$ 的“广义基础解系”)

② 非齐次线性方程组 $Ax=b$ 的任一解都可由 $\eta^*$，$\eta^*+\xi_1$，$\cdots$，$\eta^*+\xi_{n-r}$ 线性表示.\quad(解中有 $n-r(A)+1$ 个线性无关的向量)
\end{fnote}

\fsec{二、线性方程组问题}

线性方程组问题分为①一般问题；②公共解问题；③同解问题.

\fsec{1. 一般求解问题（含参常考）}

\begin{fcard}{例5.1}
设 $A$ 是秩为 2 的 $2\times5$ 矩阵，$B$ 是秩为 3 的 $3\times5$ 矩阵，若 $AB^{\mathrm T}=O$，则对于齐次线性方程组 $Ax=0$ 的任一非零解 $b$ 有（\quad）.

（A）$A^{\mathrm T}y=b$ 有唯一解\qquad\qquad（B）$A^{\mathrm T}y=b$ 有无穷多解

（C）$B^{\mathrm T}y=b$ 有唯一解\qquad\qquad（D）$B^{\mathrm T}y=b$ 有无穷多解

【解】应选（C）.

由题可得 $r(A)=2$，$r(B)=r(B^{\mathrm T})=3$.

% OPD 蓝色标注（原稿）：$D_{22}$（转换等价表述）；见到 $AB=O$，想到 $AX=O$；见到 $|A|=0$，想到 $AA^{*}=O$，$A^{*}A=O$，想到 $AX=O$，$A^{*}X=O$

又由于 $AB^{\mathrm T}=O$，即矩阵 $B$ 的每一行的转置均为齐次线性方程组 $Ax=0$ 的解，且 $Ax=0$ 的基础解系中所含解向量的个数 $s=5-2=3$. 而矩阵 $B$ 的三个行向量线性无关，故转置所对应的三个列向量是 $Ax=0$ 的一个基础解系.

又 $Ab=0$，且 $b\ne0$，于是 $b$ 可由矩阵 $B$ 的三个行向量的转置线性表示，即 $B^{\mathrm T}y=b$ 有解，且

% -----------------------------------------------------------
% PDF p64（印刷页 53）
% -----------------------------------------------------------
$r(B^{\mathrm T})=r[(B^{\mathrm T},b)]=3$，$B^{\mathrm T}y=b$ 有唯一解.

取 $A=\begin{bmatrix}1&0&0&0&0\\0&1&0&0&0\end{bmatrix}$，$B=\begin{bmatrix}0&0&1&0&0\\0&0&0&1&0\\0&0&0&0&1\end{bmatrix}$，$b=\begin{bmatrix}0\\0\\1\\0\\0\end{bmatrix}$，则 $r[(A^{\mathrm T},b)]\ne r(A^{\mathrm T})$，$r[(B^{\mathrm T},b)]=r(B^{\mathrm T})=3$，故选项（A），（B），（D）错误.
\end{fcard}

\begin{fcard}{例5.2}
设 2 阶矩阵 $A$ 满足 $a_{ij}+a_{ji}=0$，$i,j=1,2$. 实矩阵 $P=[\alpha,A\alpha]$，其中 $A\alpha$ 为非零列向量，则（\quad）.

（A）$\begin{bmatrix}A+A^{\mathrm T}&O\\E&P\end{bmatrix}x=0$ 只有零解\qquad（B）$\begin{bmatrix}A+E&O\\E&P\end{bmatrix}x=0$ 有非零解

（C）$\begin{bmatrix}O&A+A^{\mathrm T}\\P&E\end{bmatrix}x=\begin{bmatrix}\alpha\\A\alpha\end{bmatrix}$ 有无穷多解\qquad（D）$\begin{bmatrix}O&A+E\\P&E\end{bmatrix}x=\begin{bmatrix}\alpha\\A\alpha\end{bmatrix}$ 有唯一解

【解】应选（D）.

% OPD 蓝色标注（原稿）：→$D_{21}$（观察研究对象），第3讲，矩阵乘法的形状大观

由 $a_{ij}+a_{ji}=0$，知 $A^{\mathrm T}=-A$，于是 $\alpha^{\mathrm T}A\alpha=(\alpha^{\mathrm T}A\alpha)^{\mathrm T}=\alpha^{\mathrm T}A^{\mathrm T}\alpha=-\alpha^{\mathrm T}A\alpha$，即 $\alpha^{\mathrm T}A\alpha=0$，$\alpha$ 与 $A\alpha$ 正交，即 $P=[\alpha,A\alpha]$ 可逆.

对于选项（A），因为 $A^{\mathrm T}=-A$，可得 $A+A^{\mathrm T}=O$，也即
\fml{\begin{bmatrix}A+A^{\mathrm T}&O\\E&P\end{bmatrix}x=\begin{bmatrix}O&O\\E&P\end{bmatrix}x=0.}

显然系数矩阵 $\begin{bmatrix}O&O\\E&P\end{bmatrix}$ 不满秩，于是 $\begin{bmatrix}A+A^{\mathrm T}&O\\E&P\end{bmatrix}x=0$ 有无穷多解.

同理，对于选项（C），
\fml{\begin{bmatrix}O&A+A^{\mathrm T}\\P&E\end{bmatrix}x=\begin{bmatrix}O&O\\P&E\end{bmatrix}x=\begin{bmatrix}\alpha\\A\alpha\end{bmatrix},}
而又因为 $\alpha$ 为非零列向量，其增广矩阵的秩
\fml{r\begin{bmatrix}O&O&\alpha\\P&E&A\alpha\end{bmatrix}>r\begin{bmatrix}O&O\\P&E\end{bmatrix},}
可知 $\begin{bmatrix}O&A+A^{\mathrm T}\\P&E\end{bmatrix}x=\begin{bmatrix}\alpha\\A\alpha\end{bmatrix}$ 无解.

此外，由题意可知 $A=\begin{bmatrix}0&a\\-a&0\end{bmatrix}$（$a\ne0$），则 $|A+E|=\begin{vmatrix}1&a\\-a&1\end{vmatrix}=1+a^2\ne0$.

故 $\begin{vmatrix}A+E&O\\E&P\end{vmatrix}=|A+E||P|\ne0$，选项（B）只有零解，同理选项（D）有唯一解.
\end{fcard}

\begin{fcard}{例5.3}
已知矩阵 $A=\begin{bmatrix}2&1&-1\\1&-1&1\\4&5&-5\end{bmatrix}$.

（1）求一个秩为 1 的 3 阶矩阵 $C$，使得 $AC$ 为零矩阵；

% 例5.3 的「（2）」小问与两问的【解】在 PDF p65（印刷页 54），属下一分片 ch05_b；
% 本片至此结束该 fcard。
\end{fcard}
